% !PRINTED 378
\weberpage{378}
The numbers $a$, $b$, $c$ are determined by (4) only up to a common sign.

The discriminant of (3) is
\weq{6}{D = b^2 + 4ac = 4c^2 y^2 d,}
and shall also be called the \gesperrt{discriminant of the irrational number} $\omega$. $D$ is thus an integer which agrees in sign with $d$ and which is divisible by $d$; for if $p$ is a prime divisor of $d$, which therefore by hypothesis divides $d$ only once, then $p$ cannot divide the denominator of $2cy$, and hence by (6) must divide $D$. Therefore $2cy$ is also an integer, and the quotient $D:d$ is a square number.

By (4) and (6), $\omega$ and $\omega'$ can be represented thus:
\weq{7}{\omega = \frac{\sqrt{D}+b}{2c}, \quad \omega' = \frac{-\sqrt{D}+b}{2c},}
or, what is the same thing:
\weq{8}{\omega = \frac{2a}{\sqrt{D}-b}, \quad \omega' = \frac{-2a}{\sqrt{D}+b},}
if the sign of the root is determined from
\weq{9}{\sqrt{D} = 2cy\sqrt{d}}
Since an even square number is divisible by 4, while an odd one, divided by 4, leaves the residue 1, it follows from (6) that
\weq{10}{D \equiv 0 \quad \text{or} \quad \equiv 1 \ (\mathrm{mod}\ 4)}
must hold.

We wish further to investigate how the numbers $a$, $b$, $c$, $D$ change when we pass from $\omega$ to an equivalent number $\omega_1$. Let, then,
\weq{11}{\omega_1 = \frac{\alpha\omega+\beta}{\gamma\omega+\delta}, \quad \omega = \frac{\delta\omega_1-\beta}{-\gamma\omega_1+\alpha_1} \quad \cdot \quad \alpha\delta - \beta\gamma = \varepsilon.} % [sic] denominator prints $\alpha_1$, not $\alpha$; source uses a centred-dot separator, not a comma, before the third identity
From (5) one obtains for $\omega_1$ the quadratic equation
\weq{12}{c_1\omega_1^2 = a_1 + b_1\omega_1,}
if we put
\weq{13}{\begin{aligned}
-a_1 &= -a\alpha^2 + b\alpha\beta + c\beta^2\\
b_1 &= -2a\alpha\gamma + b(\alpha\delta+\beta\gamma) + 2c\beta\delta,\\
c_1 &= -a\gamma^2 + b\gamma\delta + c\delta^2
\end{aligned}}
From (13) one obtains, by solving for $a$, $b$, $c$:
%CONT>>
